% s-321-p321-m002; original PDF 321, 322, 323, 324
\Aphorism{s-321-p321-m002}{378}
\Sperr{Wiederkunft und Wiedergeburt.} Man staunt heute nicht mehr darüber, dass auf ein Zeitalter der Mathematik und der „ewigen; Pferdekräfte“ ein Zeitalter der Mystik und des Spiritualismus folgen muss; der Geschmack der Gebildeten, ungeleitet und sensationell wie er ist, kann gar nicht anders als in wüsten Extremen die Gedankenkämpfe wiedergeben, die sich in den Gehirnen der führenden Geister abspielen. Ein Naturforscher stellt redlich und vorsichtig Hypothesen auf, die nicht einmal innerhalb der Biologie durchaus hinreichend und erschöpfend sein wollen: sofort haben wir eine kühn popularisirte Entwicklungslehre als Schlüssel, mit dem der moderne Faust zu den Müttern niedersteigt In exponirten, sehr empfindlichen Intellecten finden wir einen Hang zum Unaufgehellten der menschlichen Seele, eine hyperästhetische Scheu vor dem grellen Licht natürlicher und geschichtlicher Wahrheiten: sofort fangen die kleinen Pariser Romanciers an zu katholisiren; sofort kommt Colportage-Buddhismus und Arme-Leute-Metaphysik herauf, eine unglückselige Neuauflage der alten, guten Geisterseherei, diesmal aber in scheinwissenschaftlicher Terminologie. Also nicht hier, im vergröberten Massstab populärer Neigungen, lassen sich die feinen Antriebe nachrechnen, die zu innerst ein materialistisches oder mystisches, ein bildendes oder musicirendes Zeitalter definiren; wir müssen schon zur Quelle, nämlich zu den einzelnen und temperamentstarken Individuen der Zeit, zurückgehen, um zu erfahren, was sich da begeben hat. Der Musiker des Jahrhunderts nimmt mit einem dramatischen Entwurfe von uns Abschied, der die buddhistische Wiedergeburtslehre zum Angriffspunkt der Handlung hat („Die Sieger“ von Richard Wagner); der Philosoph des Jahrhunderts nennt sich in seinen letzten Worten den Lehrer der ewigen Wiederkunft und „stellt sich damit wieder auf den Boden zurück, aus dem sein Wollen, sein Können wächst“ \Dots{4} (Nietzsche, „Götzendämmerung“). Das ist sicher ein Problem, und kein erfreuliches; denn wenn Metaphysik bei Jungen oder Verliebten eine enthusiastische Unart ist, was ist sie bei reifen und kühlen Männern, die ihre Ernte im Diesseitigen gehalten und sich der ganzen klingenden und strahlenden Weltschönheit versichert haben? Ein Undank, ein später Spieltrieb, ein Atavismus und Rückschlag ins Barbarische, oder noch anderes? Nun, vor allem eine intellectuelle Kurzsichtigkeit, ein Mangel an Schärfe und Penetranz des philosophischen Blicks. Diese beiden Lehren, Wiedergeburt und Wiederkunft, sind weder Irrthümer noch Wahrheiten, sondern Halbwahrheiten, schöne bunte Provisoria, wie sie bei jeder Sublimation philosophischer Begriffe als flüchtige Nebenproducte abfallen. Beide sind Beispiele zu dem allgemeinen Satze, dass empirische Beziehungen sich im Transcendenten nicht fixiren lassen; das empirische Ich verträgt sich mit einer transcendenten Identität der verschiedenen Ichsager, das empirische Jetzt ist transcendent vielleicht gleichzeitig mit dem empirischen Einst und Ehemals. Bleibt man auf einer Vorstufe dieser Erkenntniss stehen, so eröffnet sich allerdings die pittoreske Möglichkeit, dass der Quäler in Gestalt des Gequälten wiedergeboren, der Weltverlauf in seiner scheinbaren Unendlichkeit ewig wiederholt werde; denkt man aber den Gedanken zu Ende, so kommt man auf Allgemeinheiten von der oben ausgesprochenen Form, denen ersichtlich jede practische Nutzbarkeit, jede Ausmündung ins empirische Bewusstsein fehlt. „Ich“ werde nicht nur als die Spinne wiedergeboren, die ich zertrete, sondern als jeder beliebige Neger, den ich nie gekannt habe, als jedes Glied einer biologischen Reihe, die zu irgend einer Zeit auf irgend einem Planeten gelebt hat: die Beschränkung dieser universalen Wiedergeburt auf einen zu Büssungszwecken vorgeschriebenen Durchgang durch \Sperr{bestimmte} Lebensläufe ist Spielerei. Nicht nur die Welt als Ganzes, sondern jedes kleinste Zeittheilchen von ihr kann transcendent reproducirt werden, und in beliebiger Reihenfolge gegen andere Theilchen, aber ohne empirisches Bewusstsein davon: das heisst aber einfach, dass Vergangenheit, Gegenwart und Zukunft nur empirische Localisationen und „an sich“ nicht zu definiren sind — was man sonst aus diesem nüchternen Schulsatze macht, ist phantastischer Hocuspocus. Nietzsche hatte, wie Wagner, ein Auge für bezaubernde Halbwahrheiten, Beiden fehlt der dürre Ernst der wissenschaftlichen Verallgemeinerung.
\par

